\section{Discussion}\label{sec:discussion}

\subsection{Reading this against the coordinated-study result}\label{sec:vs_perignon}

The natural comparison is \citet{Perignon2024}, who report 52\% exact
reproduction in finance. Nothing here contradicts that number; the two measure
different situations, and the distance between them is the useful quantity.

In the coordinated setting, code is written by a participating team knowing it
will be run, the raw data are supplied, and the researchers are reachable. In
ours, the artifact is whatever a published paper points its readers to, retrieved
cold, on a clean machine, with no data supplied and nobody to ask. The first
measures whether finance code \emph{can} reproduce. The second measures whether a
reader can get it to run at all.

That the figures are 52\% and 2.8\% does not mean the coordinated result was
optimistic. It means most of what that measurement captures is supplied by
conditions surrounding the code rather than by the artifact itself: the data, the
authors' attention, the knowledge that someone will check. Strip those away and
what remains is the artifact as released, which is what a reader inherits. The
practical implication in \Cref{sec:where} follows directly: the interventions
that would move our number are the ones that make an artifact self-contained,
not the ones that make an author more careful at submission.

\subsection{What the curve licenses us to say}\label{sec:claim_scope}

The measured quantity is narrow and worth stating precisely. Of 459
quantitative-finance machine-learning papers that advertise a code repository,
thirteen executed to completion and seven produced output when a stranger cloned
the repository, built the declared environment in a container, and ran the
documented entry point without correspondence, manual repair, or reading beyond
the README.

That is not a claim that the other 452 papers are wrong, and it is not a
replication failure rate. It is a statement about artifacts, not findings: the
released code, as released, does not run unassisted on a current platform.
Several of the failures are ones an author could resolve in a message and a
motivated reader in an afternoon. The claim is about the cost of that
afternoon, multiplied across a literature.

The converse also holds, and the seven adjudicated artifacts make it concrete.
Producing output says nothing about whether the output is the paper's: two ran a
bundled self-test on synthetic inputs, one skipped every GPU measurement the
paper's claims depend on, one analyzed an empty universe after its data fetch
returned nothing, and one printed an authentication error. Clearing rung 9
establishes that a program ran, not that the study did.

\subsection{Where the attrition actually happens}\label{sec:where}

The practical value of an ordered ladder is that it localizes the loss, and the
location is not where the reproducibility discourse usually points. Link rot,
the most-discussed failure, accounts for 4.4\%. The two large drops are
environment declaration (180 repositories ship no dependency specification) and
environment construction (110 of the 252 that declare one no longer install).

These have different remedies. The first is a submission-time omission and is
cheap to fix: a \texttt{requirements.txt} clears it, and an editorial checkbox
would move the number.

The second is harder to prescribe for, because we cannot identify its mechanism.
A declared environment can fail to build because a pinned package has left the
index, because a pin set no longer admits a solution, or because a source
distribution no longer compiles, and the three are of comparable size. It is
tempting to read all of them as aging, and \Cref{tab:cohort} does show older
cohorts building far less often. But \Cref{sec:limitations} sets out why that
reading is not established: a cross-section cannot separate aging from a change
in practice, and \texttt{package\_not\_found} does not grow as a share of older
cohorts, which is what a pure index-attrition story predicts.

What can be said is weaker and still useful for policy. A journal that requires a
dependency file addresses the larger drop immediately, and that recommendation
does not depend on any mechanism. A journal that wants artifacts to survive five
years needs artifacts that pin more than package versions, in practice a
container image or a lockfile with hashes, and should know that this is a bet
against several failure modes at once rather than a fix for one identified
cause.

\subsection{Limitations}\label{sec:limitations}

\emph{The entry-point heuristic is the binding constraint on the upper rungs.}
Between rungs 6 and 8 the measurement is partly measuring us. We report both
specifications (\Cref{sec:sensitivity}) rather than choosing one, but neither
matches a determined human who would read the source, infer the intended
invocation, and supply arguments. The upper rungs are therefore a lower bound on
what a persistent reader could achieve.

\emph{The corpus is arXiv and recent.} Fields whose authors do not preprint are
absent, and the sample is weighted toward the last two years. The curve describes
current arXiv practice in quantitative finance.

\emph{The measurement is timestamped, and the direction is not established.}
Older cohorts build markedly less often (\Cref{tab:cohort}), but a single
cross-section cannot separate artifacts decaying with age from practice improving
over time, and the reason mix does not support attributing the gradient to
packages leaving the index: \texttt{package\_not\_found} is flat as a share of
dependency-declaring repositories across cohorts. What follows is only that any
comparison against these figures must be dated, and that establishing a decay
rate requires re-scoring this corpus later rather than reading aging off a
cross-section. That is part of why the harness is released rather than only the
numbers.

\emph{Time and resource budgets bind.} Dependency installation is capped at
fifteen minutes per artifact and execution at thirty, on two allocated cores and
eight gigabytes of memory
without a GPU. Four repositories hit that ceiling
and are recorded as timeouts, not failures of a different kind. An artifact whose
honest runtime is a day on a cluster is not distinguishable here from one that
hangs.

\emph{Rung 10 is human-scored, and specification-dependent.} The final rung
reflects one scorer's judgment about whether an emitted number corresponds to a
published claim. With seven eligible artifacts this is disclosed as adjudication
rather than measurement, and the worksheet, including the reasoning for every
verdict, is released so that the judgments can be checked. The count is also not
robust to the entry-point rule: the single match is reached only under the
README specification, and is zero under the conservative one
(\Cref{sec:sensitivity}). A count of one that becomes zero under a defensible
alternative should be read as a bound, not a rate. The \emph{match} was verified against the published
table cell by cell, all thirty-six values; the six \emph{not comparable} verdicts
turn on what the artifact executed rather than on any numerical tolerance, which
makes them easier to dispute and easier to check.

\subsection{The harness as the durable contribution}\label{sec:harness_contribution}

The numbers in \Cref{sec:results} describe one corpus at one moment. The
instrument that produced them is the part intended to outlast the snapshot: it is
installable, it takes any corpus of paper identifiers and code URLs, and it emits
per-artifact records with a closed failure vocabulary.

Three properties were built in because their absence would have been invisible in
the output. Harness-attributable failures carry a distinct reason code and are
excluded from claims about the literature. Conditions indicating harness
misconfiguration raise an error rather than being recorded as data, so a broken
run fails loudly instead of producing a clean-looking corpus of false negatives.
And the discretionary heuristic is a switch rather than a constant, so its effect
is reportable instead of assumed.

Those choices were not precautionary in the abstract. During development, a
container backend that installed dependencies in one container and executed in
another produced a complete, plausible, entirely false result in which nothing in
the corpus ran; it was caught because the failures took a tenth of a second. A
study that measures how often other people's code silently fails to work is
poorly placed to assume its own does.

\subsection{Future work}\label{sec:future}

Three extensions are supported by the released artifact. The corpus can be widened
beyond arXiv and beyond finance, since nothing in the ladder is domain-specific.
The same corpus can be re-scored periodically to measure the decay rate directly,
converting a limitation into a result. And the entry-point step can be
strengthened, ideally by comparing the automated heuristics against a
human-annotated subset, which would replace the sensitivity band with a
calibration.
